% ─────────────────────────────────────────────────────────────
% 封面。TikZ 绘制，不引入图片依赖，任意缩放不失真。
%
% 排版驱动，图形只做背景纹理：左下密、右上疏的节点场，
% 对应 FDE 的处境——越靠近客户现场，能依托的东西越少。
% 底部用一条实色带托住文字，从结构上杜绝纹理压字。
% ─────────────────────────────────────────────────────────────
\usepackage{tikz}
\usetikzlibrary{calc}

% ── 分册变量 ────────────────────────────────────────────────
% 两册共用这一份封面，差异只有四处文字。各册在 metadata-volN.yaml
% 的 header-includes 里 \def 覆盖；没定义时落到下面的默认值。
\providecommand{\covervolume}{}
\providecommand{\coversubtitle}{从技术栈到客户现场的\\完整落地打法}
\providecommand{\coverstats}{29 章正文\ \ ·\ \ 12 个动手 Lab}
\providecommand{\coverstatsnote}{每个 Lab 都有 verify.sh，全项通过才算完成}

\definecolor{coverbg}{HTML}{05080F}      % 近黑
\definecolor{coverband}{HTML}{0B1220}    % 底部色带，比底色略亮
\definecolor{coverink}{HTML}{F5F8FD}     % 标题白
\definecolor{coveracc}{HTML}{35D6C0}     % 唯一主色：青
\definecolor{coverdot}{HTML}{3E5A80}     % 节点色

\renewcommand{\maketitle}{%
\begin{titlepage}
\thispagestyle{empty}
\begin{tikzpicture}[remember picture, overlay]

  \fill[coverbg] (current page.south west) rectangle (current page.north east);

  % ── 节点场：左下密、右上疏 ───────────────────────────────
  \begin{scope}
    \clip (current page.south west) rectangle (current page.north east);
    \foreach \i in {0,...,17}{
      \foreach \j in {0,...,26}{
        % d：离左下角的归一化距离，越大越靠近「前沿」
        \pgfmathsetmacro{\d}{sqrt((\i/17)*(\i/17) + (\j/26)*(\j/26))/1.414}
        % 抽稀：远处成倍变疏
        \pgfmathtruncatemacro{\keep}{(\d < 0.30) ? 1 :
                                ((\d < 0.52) ? ((mod(\i+\j,2)==0) ? 1 : 0) :
                                ((\d < 0.74) ? ((mod(\i,2)==0 && mod(\j,2)==0) ? 1 : 0) :
                                               ((mod(\i,4)==0 && mod(\j,4)==0) ? 1 : 0)))}
        % 亮度与尺寸同步衰减，让梯度真的看得出来
        \pgfmathsetmacro{\op}{max(0.04, 0.62 - \d*0.74)}
        \pgfmathsetmacro{\sz}{max(0.6, 2.9 - \d*2.7)}
        \ifnum\keep=1
          \fill[coverdot, opacity=\op]
            ($(current page.south west) + (\i*1.24cm + 0.55cm, \j*1.17cm + 0.45cm)$)
            circle (\sz pt);
        \fi
      }
    }
  \end{scope}

  % ── 底部实色带：托住文字，纹理不会压上来 ─────────────────
  \fill[coverband]
    (current page.south west) rectangle
    ($(current page.south east) + (0, 5.6cm)$);
  \draw[coveracc, opacity=0.75, line width=2.4pt]
    ($(current page.south west) + (0, 5.6cm)$) --
    ($(current page.south east) + (0, 5.6cm)$);

  % ── 顶部标签 ─────────────────────────────────────────────
  \node[anchor=north west, text=coveracc, opacity=0.92,
        font=\sffamily\bfseries\fontsize{9.5}{12}\selectfont]
    at ($(current page.north west) + (2.4cm, -3.0cm)$)
    {FDE\ \ ·\ \ FORWARD  DEPLOYED  ENGINEER};
  \draw[coveracc, opacity=0.45, line width=0.7pt]
    ($(current page.north west) + (2.4cm, -3.55cm)$) --
    ($(current page.north west) + (9.6cm, -3.55cm)$);

  % ── 主标题 ───────────────────────────────────────────────
  % 断行位置就是书名的冒号：第一行是岗位，第二行是这本书的主张。
  % FDE 三个字母必须出现在封面上，而且要落在第一行——读者在书架上
  % 和招聘启事上认的是这个缩写，不是中文全称。
  %
  % 第一行 10.8 个全角字宽，41pt 下约 15.6cm，从 2.3cm 起排到 17.9cm，
  % 右边还剩 3.1cm。换更长的书名要重算这一处，不然会顶到页边。
  \node[anchor=north west, text=coverink, font=\bfseries\fontsize{41}{52}\selectfont]
    at ($(current page.north west) + (2.3cm, -5.2cm)$)
    {前线部署工程师\textcolor{coveracc}{（FDE）}};
  \node[anchor=north west, text=coverink, font=\bfseries\fontsize{41}{52}\selectfont]
    at ($(current page.north west) + (2.3cm, -7.3cm)$)
    {把 AI 交付到真实世界};

  % ── 分册徽标。没有分册时 \covervolume 为空，整块不画 ──────
  \ifx\covervolume\empty\else
    \node[anchor=north west, text=coveracc, opacity=0.95,
          font=\sffamily\bfseries\fontsize{20}{26}\selectfont]
      at ($(current page.north west) + (2.4cm, -9.7cm)$)
      {\covervolume};
  \fi

  % ── 副标题 ───────────────────────────────────────────────
  \draw[coveracc, line width=2.4pt]
    ($(current page.north west) + (2.4cm, -11.3cm)$) --
    ($(current page.north west) + (2.4cm, -13.0cm)$);
  \node[anchor=north west, text=coverink, opacity=0.86,
        font=\fontsize{14}{21}\selectfont, align=left]
    at ($(current page.north west) + (3.05cm, -11.25cm)$)
    {\coversubtitle};

  % ── 色带内：这本书的身份 ─────────────────────────────────
  \node[anchor=north west, text=coveracc, opacity=0.95,
        font=\sffamily\bfseries\fontsize{15}{19}\selectfont]
    at ($(current page.south west) + (2.4cm, 4.65cm)$)
    {\coverstats};
  \node[anchor=north west, text=coverink, opacity=0.55,
        font=\sffamily\fontsize{9.5}{13}\selectfont]
    at ($(current page.south west) + (2.4cm, 3.6cm)$)
    {\coverstatsnote};

  \draw[coverink, opacity=0.14, line width=0.6pt]
    ($(current page.south west) + (2.4cm, 2.5cm)$) --
    ($(current page.south east) + (-2.4cm, 2.5cm)$);

  \node[anchor=south west, text=coverink, opacity=0.9, font=\fontsize{12}{16}\selectfont]
    at ($(current page.south west) + (2.4cm, 1.45cm)$)
    {whiteyoungy};
  \node[anchor=south east, text=coveracc, opacity=0.8,
        font=\sffamily\fontsize{12}{16}\selectfont]
    at ($(current page.south east) + (-2.4cm, 1.45cm)$)
    {2026};
  \node[anchor=south west, text=coverink, opacity=0.45,
        font=\sffamily\fontsize{8.5}{11}\selectfont]
    at ($(current page.south west) + (2.4cm, 0.75cm)$)
    {© 2026 whiteyoungy　·　非商业使用依 CC BY-NC 4.0 开放　·　商业使用须书面授权：whiteyoung@gmail.com};

\end{tikzpicture}
\end{titlepage}
}
